%% Establishing the phenomenon on models trained without any OPERA mechanism (Tables 2-3).
%% Included by sigconf.tex. Do not add a preamble here.

\section{Does Observation-Process Shift Matter?}
\label{sec:phenomenon}
We first establish the phenomenon on models trained without any OPERA mechanism.

\paragraph{Severity spans three orders of magnitude.}
Table~\ref{tab:phenomenon} reports each baseline under its own observation process and under shift. Mean
degradation ranges from $+1.9\%$ to $+1793\%$, and the ordering is interpretable: it grows with how much of the
task's predictive signal must come from the irregular history. USHCN forecasting one month ahead from five
climate variables sits at the bottom, and the diagnostic is unambiguous --- removing half of its observations at
matched density \emph{lowers} its error, from $0.5379$ to $0.5355$. Its signal is carried by seasonal structure
that survives the loss. Extending the same dataset to a twelve-month horizon doubles its sensitivity
($+1.9\%\to+3.7\%$), which identifies the mechanism of the null rather than excusing it. At the top, Activity's
millisecond-resolution sensor dynamics degrade by $32\times$ under stride-4 subsampling alone.

\begin{table}[t]
\caption{Observation-process shift on models trained without any OPERA mechanism\prelim. Primary metric per
setting; degradation by Eq.~\ref{eq:degradation}. Every shifted condition re-observes only the history.}
\label{tab:phenomenon}
\small
\begin{tabular}{lrrrr}
\toprule
setting & clean & worst condition & $\Delta_{\mathrm{shift}}$ & $\Delta_{\mathrm{unseen}}$ \\
\midrule
Activity (MSE) & 0.00263 & 0.08547 \textit{(stride 4)} & $+1793\%$ & $+2057\%$ \\
PhysioNet (MSE) & 0.00549 & 0.02134 \textit{(var .5)} & $+107\%$ & $+130\%$ \\
PAM (acc.) & 88.01 & 53.93 \textit{(var .5)} & $+64\%$ & $+88\%$ \\
USHCN-12mo (MSE) & 0.52020 & 0.55338 \textit{(comp.)} & $+3.7\%$ & $+4.6\%$ \\
P19 (AUROC) & 82.13 & 74.63 \textit{(var .3)} & $+2.9\%$ & $+3.3\%$ \\
USHCN-1mo (MSE) & 0.53787 & 0.56933 \textit{(var .5)} & $+1.9\%$ & $+2.8\%$ \\
P12 (AUROC) & 83.53 & 77.98 \textit{(comp.)} & \multicolumn{2}{c}{\textit{see \S\ref{sec:analysis}}} \\
\bottomrule
\end{tabular}
\end{table}

\paragraph{Density does not explain it.}
Table~\ref{tab:density} compares processes that remove approximately half of the observations but destroy
different structure. On PhysioNet the spread is $2.9\times$ and on PAM $3.5\times$ in cross-entropy; a model can
lose half its observations almost for free if they are removed uniformly in time, and catastrophically if they
are removed by variable. Treating irregularity as a scalar corruption level cannot represent this distinction,
which is the empirical case for modelling the observation process as an object.

\begin{table}[t]
\caption{Matched-density comparison\prelim: four processes each removing $\approx$50\% of observations, on
models trained without OPERA. Error varies by $2.9$--$3.5\times$ with \emph{what} is removed.}
\label{tab:density}
\small
\begin{tabular}{lrrrr}
\toprule
 & density- & stride & burst & variable \\
setting & matched & 2 & .5 & dropout .5 \\
\midrule
PhysioNet (MSE) & 0.00736 & 0.00791 & 0.00804 & \textbf{0.02134} \\
Activity (MSE) & 0.04892 & 0.05230 & \textbf{0.01973} & 0.05778 \\
PAM (cross-ent.) & 0.4466 & 0.4360 & 0.4897 & \textbf{1.5578} \\
\bottomrule
\end{tabular}
\end{table}

\paragraph{The damaging kind of loss is modality-dependent.}
On the clinical datasets, losing whole variables dominates while temporal subsampling is nearly free: on P19,
removing 30\% of variables costs 7.5 AUROC whereas stride 3 costs 1.1; on PAM, removing half the variables costs
34 accuracy points whereas stride 3 costs 1.5. Activity inverts this, with stride its single worst condition and
burst removal its mildest (Table~\ref{tab:density}). Hourly laboratory values are largely redundant in time and
irreplaceable across channels; body-sensor dynamics are the reverse. A robustness method must therefore cover
process \emph{families}, not a severity axis.
